\begin{cabstract}
本报告是修读学科技术前沿讲座课程后的综合总结。课程邀请了北京航空航天大学计算机学院多位学者，围绕计算机科学的核心问题、科研方法以及多个前沿技术领域进行深入讲解。内容涵盖计算机科学哲学与科研工作方法、虚拟现实与增强现实及数字孪生技术、人工智能与大模型（以DeepSeek为例）、人形机器人与具身智能、高性能计算、多维度视觉计算以及存算传融合的自主计算架构。

本文首先对各讲座的核心内容进行系统概括，梳理从学科通论到关键技术再到系统融合的知识脉络。然后重点阐述个人在学习过程中的深度思考与收获，包括对科研本质的重新认识、对人工智能热潮的理性分析、对技术范式转换的理解以及对个人未来学习和研究方向的启发。报告最后对课程进行总结，并展望智能化时代背景下计算机学子应有的定位与追求。

本报告体现了个人独立的分析与思考。
\end{cabstract}

\begin{eabstract}
This report is a comprehensive summary after taking the course Frontier Lectures on Disciplinary Technology. Several renowned scholars from the School of Computer Science at Beihang University gave in-depth lectures on core issues of computer science, research methodologies, and multiple cutting-edge technology fields. The contents cover philosophy of computer science and research methods, virtual reality and augmented reality and digital twins, artificial intelligence and large language models (taking DeepSeek as an example), humanoid robots and embodied intelligence, high-performance computing, multi-dimensional visual computing, and the autonomous computing paradigm integrating storage, computing, and networking.

This report first systematically summarizes the core content of each lecture, outlining a knowledge context from general theory to key technologies and then to system integration. On this basis, it focuses on my in-depth thinking and gains during the learning process, including a renewed understanding of the essence of scientific research, a rational analysis of the AI boom, a deep insight into the paradigm shift from modularization to integration, and inspirations for my future study and research directions. Finally, the report concludes the course and looks forward to the positioning and pursuit that computer science students should have in the intelligent era.

This report reflects my independent analysis and thinking.
\end{eabstract}

% 关键词
\ckeyword{学科前沿，人工智能，大模型，虚拟现实，具身智能，科研方法}
\ekeyword{Frontiers of Discipline, Artificial Intelligence, Large Model, Virtual Reality, Embodied Intelligence, Research Methodology}